\section{Coupling the two Weyl blocks}\label{sec:coupling}

Assume that neither Weyl block vanishes identically and that their second
moments agree:
\[
 \mathbf E v^2=\mathbf E b^2.
\]
We retain the notation
\[
 \delta=2-v-b,\qquad K=4+v^2+b^2,\qquad
 M=1-\frac38(v-b),\qquad m=M^2,
\]
and place a bar over an expression to exchange the two blocks. In
particular, $\bar M=1-\frac38(b-v)$ and $\bar m=\bar M^2$.

The weights in \eqref{eq:12} have opposite slopes along $v+b=2$.
Their derivatives introduce quartic terms in the first derivatives of
curvature. We combine the two Hessian inequalities with one identity that
controls those costs, the balanced variance defect, and tests of the weak
norm equations. This produces a single integrand with nonnegative mean.
The task is then to bound its quartic, quadratic, and potential parts by
three expressions that complete one square.

\subsection{The combined integral inequality}

Start with the sum of the two integrands in \eqref{eq:10}. The three
additions below are written explicitly in \eqref{eq:16}--\eqref{eq:18}.
Here $h\ge0$ is a scalar coefficient and $\kappa$ is a scalar multiplier,
whose rational values are specified below. Their roles and integral signs
are as follows; the degrees refer to first
derivatives of curvature in the resulting integrand.
\begin{center}
\small
\begin{tabular}{@{}p{0.24\textwidth}p{0.42\textwidth}p{0.25\textwidth}@{}}
\toprule
Added expression & Integral and derivative degrees & Purpose\\
\midrule
$\Delta\delta$ identity &
$\dfrac h3\mathbf E[\delta(|Y|^2+|\bar Y|^2)]\ge0$; degrees $4,2,0$ &
Controls the quartic weight costs\\
Balanced variance defect & Nonnegative integral; degrees $2,0$ &
Controls $v-b$ using moment balance\\
$\kappa$ norm tests & Zero integral; degrees $2,0$ &
Adjusts the quadratic and potential terms\\
\bottomrule
\end{tabular}
\end{center}

We use the following verified rational choice for the constants and
functions in this combination and its final comparison. It is one
sufficient choice; no optimality or uniqueness is asserted.
Set
\[
 D=158,\qquad h=\frac{17}{10},\qquad
 c=\frac25-\frac3{10}vb,
 \qquad a=v-1,\quad s=b-1,\quad
 n=\frac\delta2,\quad t=\frac{v-b}{2}.
\]
Thus $c\ge 1/10$ on $v,b\ge0$, $v+b\le2$. Define
\begin{equation}
\begin{aligned}
 \kappa(v,b)={}&-\frac D6(a-s)+z_1a^2
       +\left(60-\frac{2D}{9}+2z_1\right)as+z_2s^2\\
 &+\sum_{i=0}^3z_{3+i}a^is^{3-i}
       +\sum_{i=0}^4z_{7+i}a^is^{4-i},\\
 4(z_1,\ldots,z_{11})={}&
       (-43,-45,-63,21,183,208,61,46,101,79,-41),\\
 4\Theta(v,b)={}&n^2(1-n)(105+5n)
       +t^2(4+360n-586n^2)-4t^4,\\
 R_v={}&\frac{4(1+v)}{1+6v+v^2},\qquad
 R_b=\frac{4(1+b)}{1+6b+b^2}.
\end{aligned}
\tag{15}\label{eq:15}
\end{equation}
Here and below an exchanged expression is obtained by exchanging all
block variables, and $\bar\kappa=\kappa(b,v)$.
The multiplier $\kappa$ changes the integrand through a zero-integral
norm test. The functions $\Theta$, $c$, $R_v$, and $R_b$ enter only the
comparison for the final square; they are not additional integrands.
Neither $\kappa$ nor $\Theta$ is required to have a sign.

The first addition is
\begin{equation}
\begin{aligned}
 h\bigg[&(e+\bar e)
   \left(\lambda-\frac{e-|p|^2}{v}
                     -\frac{\bar e-|q|^2}{b}\right)
       +40\delta(B+\bar B)\\
 &\hspace{24mm}
       -30\delta\sum_{\pm}v^2(1-2\sigma v+v^2)\bigg],
\end{aligned}
\tag{16}\label{eq:16}
\end{equation}
where
\[
 \lambda=6(\sigma v^2+\bar\sigma b^2-v-b)-26\delta.
\]
The notation $\sum_\pm$ includes the displayed term and its exchange.
To verify its sign, use
\[
 \Delta e=\frac{2}{6}|\nabla Z|^2+26e-40B
\]
and \eqref{eq:3} in
\[
 \mathbf E\bigl[(e+\bar e)\Delta\delta\bigr]
 =\mathbf E\bigl[\delta\Delta(e+\bar e)\bigr].
\]
The decomposition $\nabla Z=Y+D_0$ gives
$|D_0|^2=(5/12)|\Delta T|^2$. Moreover,
\[
 \frac{2|D_0|^2}{6}
       =30v^2(1-2\sigma v+v^2).
\]
Consequently the integral of \eqref{eq:16} equals
$\frac h3\mathbf E[\delta(|Y|^2+|\bar Y|^2)]$, which is nonnegative.

The second addition is
\begin{equation}
 D\left[\frac{|p-q|^2}{K}
       -\left(1-\frac{\delta^2}{4}\right)(v-b)^2\right].
\tag{17}\label{eq:17}
\end{equation}
Indeed, the equality of the second moments implies
\[
 \mathbf E(v-b)=\frac12\mathbf E[\delta(v-b)],
\]
and hence
\[
 (\mathbf E(v-b))^2\le\frac14\mathbf E[\delta^2(v-b)^2].
\]
Applying \eqref{eq:4} to $v-b$ proves that \eqref{eq:17} has
nonnegative integral. Finally add the zero-integral expression
\begin{equation}
 \sum_\pm\left[
 \kappa(v,b)\left\{6v(1-\sigma v)+\frac{e-|p|^2}{v}\right\}
       +\nabla\kappa(v,b)\cdot p\right].
\tag{18}\label{eq:18}
\end{equation}
This is the weak identity \eqref{eq:3} tested against $\kappa(v,b)$.

Let $\mathcal I$ denote the resulting total integrand. Thus
$\mathbf E\mathcal I\ge0$. All these integrals are legitimate without
a positive lower bound for either norm: $e/v$ and $\bar e/b$ are
integrable by the regularization proving \eqref{eq:3}, while the energies
and all scalar coefficients are bounded. Bounded Lipschitz test functions
are admissible by smooth approximation, uniformly and in $H^1$.

For the fixed weights in \eqref{eq:12}, $dm$ is linear in $p,q$.
Thus the cubic contraction multiplied by $dm$ and the quadratic form
in $L_m=dm\otimes Z$ have degree four in first derivatives of curvature.
Grouping the exact total integrand by this degree gives
$\mathcal I=\mathcal I_4+\mathcal I_2+\mathcal I_0$, where
\[
\begin{aligned}
\mathcal I_4={}&
 \sum_\pm\left[
 -\frac56\sum_i m_i\,
   \frac{\operatorname{tr}\bigl(Z_i\sum_j Z_j^2\bigr)}6
 +\frac{\langle\mathcal P(L_m)_Y,(L_m)_Y\rangle}{24m}
 \right]\\
 &-h(e+\bar e)
 \left(\frac{e-|p|^2}{v}+\frac{\bar e-|q|^2}{b}\right),\\[1mm]
\mathcal I_2={}&
 \sum_\pm\left[
 mG_v+h\lambda e+40h\delta B
 +\kappa(v,b)\frac{e-|p|^2}{v}
 +\nabla\kappa(v,b)\cdot p
 \right]
 +\frac DK|p-q|^2,\\[1mm]
\mathcal I_0={}&
 -D\left(1-\frac{\delta^2}{4}\right)(v-b)^2\\
 &+\sum_\pm\left[
 6\kappa(v,b)v(1-\sigma v)
 -30h\delta v^2(1-2\sigma v+v^2)
 \right].
\end{aligned}
\]
Here $\lambda$ is the original symmetric expression in \eqref{eq:16}.
The exchange in $\sum_\pm$ includes the operator, the $Y$-projection,
the weight, the multiplier and its arguments, and all derivative labels.
On a block's zero set, only its corresponding derivative summands vanish
almost everywhere; the exchanged contributions remain. The quotients
use the zero value assigned in \eqref{eq:3}, and $G_v$ uses its invariant
expression there. Products such as $\sigma v^2$ use their continuous
zero value; no continuity is asserted for the singular coefficients.

\begin{samepage}
\begin{proposition}[Three bounds for the combined integrand]
\label{prop:coupling-estimates}
For this decomposition, define
\[
 r=R_v|p|^2+R_b|q|^2,\qquad
 j=(|p|^2+|q|^2)
       \left(\frac{|p|^2}{v}+\frac{|q|^2}{b}\right).
\]
Then, almost everywhere,
\[
 \mathcal I_4\le-cj,\qquad
 \mathcal I_2\le2\Theta r,\qquad
 \mathcal I_0\le-\frac{\Theta^2}{c}.
\]
The corresponding derivative summands are assigned value zero on a
zero set, as in \eqref{eq:3}. The total integrand and the bounds are
integrable, and $\mathbf E\mathcal I\ge0$.
\end{proposition}
\end{samepage}

Section~\ref{sec:classification} combines these estimates with $r^2\le j$ to obtain
\[
 \mathcal I\le-\frac{(\Theta-cr)^2}{c}-c(j-r^2).
\]
The strictness needed for vanishing is proved there.
We prove the three estimates in turn. The computations first take place
on $v,b>0$, $v+b\le2$, $x+y\le2$. The polynomial certificates used in
the proof are supplied in Section~\ref{sec:certificates}.

\subsection{The quartic terms}

Put
\[
 e_0=\frac53|p|^2,\qquad e_d=e-e_0\ge0,\qquad z=|p|+|q|.
\]
Since $dm=-\frac34M(p-q)$, \eqref{eq:13} bounds the cubic contraction
multiplied by $dm$ in \eqref{eq:10} by
$\frac12\sqrt{3/5}\,Mz e^{3/2}$. The nonnegative bracket in
\eqref{eq:14} is bounded above by $g e_0+H_d e_d$, where
\[
 g=\frac7{12}(4+y)-\sigma v+\frac{\rho^2v^2}{4},\qquad
 H_d=\frac7{12}(4+y)+\frac54v+4.
\]
In fact, the pure defect contribution to $-5B/4$ is at most
$(5/4)v e_d$, and its mixed contribution is at most
$2v|\rho|\sqrt{e_0e_d}$. The latter is at most
$(\rho^2v^2/4)e_0+4e_d$. Nonnegativity of the original bracket permits
the replacement $|p-q|^2\le z^2$. We obtain
\begin{equation}
\begin{aligned}
 \mathcal I_4\le{}&
 \sum_\pm\left[
       \frac12\sqrt{\frac35}\,zMe^{3/2}
       +\frac9{64}z^2(g e_0+H_de_d)\right]\\
 &-h(e+\bar e)
       \left[\frac{(2/5)e_0+e_d}{v}
             +\frac{(2/5)\bar e_0+\bar e_d}{b}\right].
\end{aligned}
\tag{19}\label{eq:19}
\end{equation}

We first show that this upper bound decreases when either defect
increases, with $p,q$ fixed. Denote the right-hand side by $\mathcal Q$.
For $e>0$, direct differentiation gives
\[
\begin{aligned}
 \partial_{e_d}\mathcal Q={}&
       \frac34\sqrt{\frac35}\,zM\sqrt e
       +\frac9{64}z^2H_d\\
 &-h\left[
       \frac{(2/5)e_0+e_d}{v}
       +\frac{(2/5)\bar e_0+\bar e_d}{b}
       +\frac{e+\bar e}{v}\right].
\end{aligned}
\]
Set $u=\sqrt{\bar e/e}$. The bounds
\[
\begin{gathered}
 z\le\sqrt{\frac35}(\sqrt e+\sqrt{\bar e}),\qquad
 M\le\frac74-\frac34v,\qquad H_d\le\frac{15}{2}+\frac23v,\\
 \frac25e_0+e_d\ge\frac25e,\qquad
 \frac25\bar e_0+\bar e_d\ge\frac25\bar e,
 \qquad \frac vb\ge\frac{v}{2-v}
\end{gathered}
\]
therefore imply
\[
 \frac ve\partial_{e_d}\mathcal Q
 \le A_0(1+u)^2+L_0(1+u)
       -h\left[\frac75+
          \left(1+\frac{2v}{5(2-v)}\right)u^2\right],
\]
where
\[
 A_0=\frac{27}{320}v\left(\frac{15}{2}+\frac23v\right),
 \qquad L_0=\frac9{20}v\left(\frac74-\frac34v\right).
\]
All substitutions in the negative bracket use lower bounds; thus they
give the required upper bound after multiplication by $-h$.

Write
\[
 C_0=\frac75h-A_0-L_0,\qquad
 D_0^{\mathrm{sc}}=h-A_0+\frac{2hv}{5(2-v)}.
\]
Lemma~\ref{lem:scalar-certificates} gives $C_0>0$ and
\[
 4C_0\left((h-A_0)(2-v)+\frac25hv\right)
       -(2-v)(2A_0+L_0)^2>0.
\]
For $0<v<2$ this says
$4C_0D_0^{\mathrm{sc}}>(2A_0+L_0)^2$; hence the homogenized
quadratic
\[
 -C_0+(2A_0+L_0)u-D_0^{\mathrm{sc}}u^2
\]
is negative for every real $u$. No ordering of the two energies is
needed. Before taking an endpoint limit, write the derivative bound as
\[
 \partial_{e_d}\mathcal Q\le\frac1v
 \left[-C_0e+(2A_0+L_0)\sqrt{e\bar e}
                    -D_0^{\mathrm{sc}}\bar e\right].
\]
It extends continuously to $e=0$, where necessarily $e_0=e_d=0$.
At $e=\bar e=0$ its value is zero; otherwise the derivative remains
negative. The exchanged argument applies to $\bar e_d$. Decreasing
the defects one at a time consequently gives
\[
 \mathcal Q(e_0+e_d,\bar e_0+\bar e_d)
 \le\mathcal Q(e_0,\bar e_0+\bar e_d)
 \le\mathcal Q(e_0,\bar e_0).
\]

We may therefore put both defects equal to zero. Write
\[
 V=x-v,\qquad W=y-b,\qquad
 0\le V\le v,\quad0\le W\le b,\quad V+W\le\delta.
\]
We claim
\[
 g\le g_v:=\frac7{12}(4+b+W)-v+2V,
 \qquad
 g_b=\frac7{12}(4+v+V)-b+2W.
\]
For $r_0=x/v\in[1,2]$, use
\[
 \sigma=\frac{3r_0-r_0^3}{2},\quad
 1-\sigma=\frac{(r_0-1)^2(r_0+2)}2,\quad
 1+\sigma=\frac{(2-r_0)(r_0+1)^2}{2}.
\]
Since $r_0v\le2$,
\[
 \frac{(r_0-1)(r_0+2)}2
       \left(1+\frac{v(1+\sigma)}4\right)\le2.
\]
After substituting $v\le2/r_0$, the gap in this inequality is
\[
 \begin{aligned}
 &2-\frac{(r_0-1)(r_0+2)}2
       \left(1+\frac{1+\sigma}{2r_0}\right)\\
 &\quad=\frac{2-r_0}{8r_0}
 \left[4r_0(r_0+3)-(r_0-1)(r_0+2)(r_0+1)^2\right]\ge0.
 \end{aligned}
\]
The bracket is concave on $[1,2]$ and has endpoint values $16$ and $4$.
This also covers $r_0=2$ without division by $2-r_0$. Multiplying the preceding bound
by $V=v(r_0-1)$ proves the claim for $g$, using
$\rho^2=1-\sigma^2$.

Set $P=\sqrt{e_0}$ and $Q=\sqrt{\bar e_0}$. We will prove
\begin{equation}
\begin{aligned}
 &\left(\frac{67}{50}+\frac{27}{100}vb\right)
       (P^2+Q^2)(bP^2+vQ^2)\\
 &\qquad\ge
       \frac{27}{128}vb(P+Q)^2(g_vP^2+g_bQ^2)
       +\frac34vb(P+Q)(MP^3+\bar M Q^3).
\end{aligned}
\tag{20}\label{eq:20}
\end{equation}
At zero defects, $z=\sqrt{3/5}(P+Q)$. Thus the positive terms of
\eqref{eq:19} are at most $2/5$ times the right-hand side of
\eqref{eq:20}, divided by $vb$. Since
\[
 \frac{10}{9}\left(h-\frac{67}{50}-\frac{27}{100}vb\right)=c,
\]
\eqref{eq:20} gives the first estimate in
Proposition~\ref{prop:coupling-estimates}; the factor $9/25$ that remains
is exactly the conversion from $(e_0,\bar e_0)$ to $(|p|^2,|q|^2)$.

To prove \eqref{eq:20}, exchange the blocks, if necessary, so that
$P>0$ and $0\le u=Q/P\le1$; the case $P=Q=0$ is immediate.
In $g_v+g_bu^2$ the coefficients of $V,W$ are nonnegative, and their
difference is $(17/12)(1-u^2)\ge0$. Replacing $(V,W)$ by
$(\delta,0)$ therefore enlarges this expression. This is a relaxation,
not an assertion that the new pair is geometrically attained. The
enlarged values are
\[
 g_v=\frac{19}{3}-3v-\frac{17}{12}b,
 \qquad g_b=\frac72-\frac{19}{12}b.
\]
Put $\omega=27/128$, $\mu=9/32$, and define
\[
\begin{aligned}
 A&=\frac{67}{50}(1+u^2),\\
 G&=\frac{27}{100}(1+u^2)u^2
       +3\omega(1+u)^2+\mu(1+u)(1-u^3),\\
 H&=\frac{27}{100}(1+u^2)
       +\frac\omega{12}(17+19u^2)(1+u)^2
       -\mu(1+u)(1-u^3),\\
 L&=\omega(1+u)^2\left(\frac{19}{3}+\frac72u^2\right)
       +\frac34(1+u)(1+u^3).
\end{aligned}
\]
Dividing the difference between the left and enlarged right sides of
\eqref{eq:20} by $vbP^4$ gives exactly
\[
 \frac Av+\frac{Au^2}{b}+Gv+Hb-L.
\]
Here $A,L>0$, and Lemma~\ref{lem:scalar-certificates} gives $G,H>0$.
For
\[
 Z_*=L^2-4A(G+u^2H)
\]
the same lemma gives $Z_*\le0$ on $[0,1/16]$ and
$Z_*^2\le64A^2u^2GH$ on $[1/16,1]$. In the latter interval,
$Z_*\le|Z_*|\le8Au\sqrt{GH}$. In both intervals it follows that
\[
 L\le2\sqrt A(\sqrt G+u\sqrt H).
\]
Taking square roots is legitimate because both sides are nonnegative.
The arithmetic--geometric mean inequalities applied to $A/v+Gv$ and
$Au^2/b+Hb$ finish the proof of \eqref{eq:20}.

\subsection{The quadratic terms}

The quartic bound is established. We next bound the quadratic part by its \emph{axial value}: the algebraic
comparison with $\sigma=\bar\sigma=1$, $\rho=\bar\rho=0$, and
$k=l=\bar k=\bar l=0$, holding $v,b,p,q$ fixed. The corresponding spectral
shapes are $(-v,-v,2v)$ and $(-b,-b,2b)$. In the plus contribution, replacing
$\bar\sigma$ by $1$ raises the expression by
$6h(1-\bar\sigma)b^2e\ge0$. After this replacement, the shape-dependent
quadratic expression is
\[
 mG_v+h\lambda_v e+40h\delta B
 +\frac{\kappa}{v}(e-|p|^2)
 +(\partial_v\kappa)|p|^2+(\partial_b\kappa)p\cdot q,
\]
where $\lambda_v=6(\sigma v^2+b^2-v-b)-26\delta$. The last two terms and
the separate variance term are unchanged by the shape comparison.
Define
\[
 P_0=300m+30hv+160h\delta,\qquad
 P_k=300m+160h\delta,\qquad P_l=35mv.
\]
On the closed spectral triangle, $m\ge1/16$ and $h,\delta,v\ge0$;
therefore $P_0=P_k+30hv>0$. This supplies the positive denominator in
the Schur comparison below.
By \eqref{eq:11} and \eqref{eq:16}, the loss in the $pp$ entry when
$\sigma<1$ is $(1-\sigma)vP_0/3$, and the mixed entries with $(k,l)$
are $-v\rho(P_k,P_l)/\sqrt3$. The negative of the pure $(k,l)$ matrix
is exactly
\[
 \mathsf D_\sigma=
 \begin{pmatrix}K_1&K_{12}\\K_{12}&K_2\end{pmatrix}
       -\frac\kappa v\operatorname{diag}(5,1),
\]
where
\[
\begin{aligned}
 K_2&=m(42+3b^2+10v^2)-h\lambda_v,\\
 K_1&=5K_2+m(300v\sigma+140v^2)+160h\delta v\sigma,\\
 K_{12}&=m(75v+35v^2\sigma)+40h\delta v.
\end{aligned}
\]
Thus the required completion of squares follows from
\begin{equation}
 P_0\left[
 v\begin{pmatrix}K_1&K_{12}\\K_{12}&K_2\end{pmatrix}
       -\kappa\operatorname{diag}(5,1)\right]
 -(1+\sigma)v^2
       \begin{pmatrix}P_k\\P_l\end{pmatrix}
       \begin{pmatrix}P_k&P_l\end{pmatrix}\succeq0.
\tag{21}\label{eq:21}
\end{equation}

Let $\mathsf A_\sigma$ be this matrix and $w=(P_k,P_l)^T$.
The identity $P_0=P_k+30hv$ gives
\[
 -\partial_\sigma\mathsf A_\sigma
 =v^2\left[
       \begin{pmatrix}30hv\\-P_l\end{pmatrix}
       \begin{pmatrix}30hv&-P_l\end{pmatrix}
       +\operatorname{diag}(0,6hvP_0)\right]\succeq0.
\]
Lemma~\ref{lem:matrix-certificates} proves \eqref{eq:21} at
$\sigma=1$, and this matrix monotonicity proves it for every
$-1\le\sigma\le1$. When $\sigma<1$, put
$\ell=(1-\sigma)vP_0/3>0$ and $d=-v\rho w/\sqrt3$.
The Schur complement is
\[
 \mathsf D_\sigma-\frac{dd^T}{\ell}
 =\mathsf D_\sigma-\frac{(1+\sigma)v}{P_0}ww^T
 =\frac{\mathsf A_\sigma}{P_0v}\succeq0.
\]
At $\sigma=1$, one has $\rho=0$ and
\[
 \mathsf D_1=\frac{\mathsf A_1+2v^2ww^T}{P_0v}\succeq0,
\]
so no division by $1-\sigma$ is needed. No sign of $\kappa$ is assumed.
Applying the exchanged argument proves the announced axial reduction.

At the axial comparison, \eqref{eq:2} and the matrix for $B/v$ give
$e_{\rm ax}=5|p|^2/3$ and $B_{\rm ax}=4v|p|^2/3$. Put
$\lambda_{\rm ax}=6(v^2+b^2-v-b)-26\delta$. Since $\delta=-a-s$,
\[
\begin{aligned}
\bigl(h\lambda e+40h\delta B\bigr)_{\rm ax}
 &=\frac53h\bigl(\lambda_{\rm ax}+32v\delta\bigr)|p|^2\\
 &=\frac53h\bigl(6(a^2+s^2)+32a\delta\bigr)|p|^2.
\end{aligned}
\]
Together with the axial $pp$ entry of \eqref{eq:11}, this gives
$X_v|p|^2$ for the axial part of $mG_v+h\lambda e+40h\delta B$, where
\[
 X_v=m(-70+100v-25v^2-5b^2)
       +\frac53h\bigl(6(a^2+s^2)+32a\delta\bigr),
 \qquad X_b=X_v(b,v).
\]
Collecting the remaining quadratic terms gives the desired estimate
precisely when
\begin{equation}
\begin{gathered}
 \begin{pmatrix}\mathcal Q_v&J\\J&\mathcal Q_b\end{pmatrix}
       \preceq2\Theta\operatorname{diag}(R_v,R_b),\\
 \mathcal Q_v=X_v+\frac DK+\partial_v\kappa+\frac{2\kappa}{3v},
 \qquad
 \mathcal Q_b=X_b+\frac DK+\partial_b\bar\kappa
                                      +\frac{2\bar\kappa}{3b},\\
 J=\frac{\partial_b\kappa+\partial_v\bar\kappa}{2}-\frac DK.
\end{gathered}
\tag{22}\label{eq:22}
\end{equation}
The polynomial matrix certificate in Lemma~\ref{lem:matrix-certificates}
will prove this inequality.

The leading terms of $\kappa$ can now be read in their intended role.
Along the formal axial boundary $v=1+t$, $b=1-t$, direct substitution
in the displayed formulas gives
\[
 \mathcal Q_v,\mathcal Q_b,J=O(t^2)\quad(t\to0),
 \qquad \Theta=t^2-t^4.
\]
Thus the constant and tangential linear terms cancel at the meeting
point of the two axial blocks. This explains the local shape of the
rational choice; the certificate proves its inequality on the full
stated domain.

\subsection{The potential and the zero sets}

The quadratic estimate has been reduced to its explicit matrix certificate.
It remains to bound the derivative-free potential and to extend the
three inequalities across the zero sets. The potential is symmetric under exchanging the two blocks. For its
pointwise estimate, suppose $v\ge b$ and define
\[
\begin{aligned}
 E_0={}&D\left(1-\frac{\delta^2}{4}\right)(v-b)^2
       +30h\delta(v^2a^2+b^2s^2)
       +6(av\kappa+sb\bar\kappa),\\
 F={}&cE_0-\Theta^2,\qquad
 L_v=10h\delta v-\kappa,\qquad
 L_b=10h\delta b-\bar\kappa.
\end{aligned}
\]
Direct collection from \eqref{eq:16}--\eqref{eq:18} gives
\[
 \mathcal I_0=-E_0-6v^2(1-\sigma)L_v
                    -6b^2(1-\bar\sigma)L_b.
\]
Since $y-b\le\delta$ and $r_0=y/b\le2$,
\[
 b^2(1-\bar\sigma)
 =\frac12b^2(r_0-1)^2(r_0+2)\le2\delta^2.
\]
It therefore suffices to establish
\begin{equation}
 L_v\ge0,\qquad F\ge0,\qquad
 F+12c\delta^2L_b\ge0\qquad(v\ge b).
\tag{23}\label{eq:23}
\end{equation}
Indeed, if $L_b\ge0$ then $\mathcal I_0\le-E_0\le-\Theta^2/c$.
If $L_b<0$, the displayed spectral bound instead gives
\[
 \mathcal I_0\le-E_0-12\delta^2L_b\le-\Theta^2/c.
\]
Lemmas~\ref{lem:matrix-certificates} and
\ref{lem:potential-certificate} prove all three inequalities in
\eqref{eq:23}.

The use of an ordering in \eqref{eq:22} and \eqref{eq:23} is purely
algebraic. All weights and their derivatives have already been formed
from the fixed global labels of the two blocks. Their summed targets
are symmetric. At a point with $v<b$, apply the ordered inequality to
$(b,v)$ and exchange the two vector entries. At $v=b$, either labeling
gives the same conclusion. No maximum, minimum, or selected label is
differentiated.

Finally, on the actual zero set $\{v=0\}$, one has $x=0$ and $y<2$ by
strict sectional positivity, and the components $Z,p,k,l$ vanish almost everywhere.
To pass the algebraic bounds to this set, first set all plus derivative
parameters $p,k,l$ to zero in the positive-variable inequalities.
Then $e,B,C_N$ and every numerator divided by $v$ are identically zero
before taking any limit. Keep $b>0$ and all minus spectral and derivative
data fixed. Choose axial plus spectra $(-v,-v,2v)$, for which $x=v$,
with $0<v<\min\{2-y,2-b\}$ tending to zero. This sequence remains in
both required spectral inequalities. The remaining quadratic inequality
is the minus diagonal, whose denominator $bK(1+6b+b^2)$ stays positive;
the quartic bound becomes $-c|q|^4/b$. All cross terms with $p$ vanish.
This argument does not assert continuity of the singular plus
coefficients. The other zero set is treated by
exchange, and if both norms vanish all derivative components vanish.
The potential inequalities are polynomial and extend continuously.
The quartic quotients are bounded in absolute value by a constant times
$e/v+\bar e/b$, and are therefore integrable. This completes the proof
of Proposition~\ref{prop:coupling-estimates}, subject only to the explicit
polynomial verifications in the next section.
